\subsection{ChucK}\label{subsec:chuck}

\textbf{ChucK} is a strongly-timed, concurrent programming language for real-time audio synthesis, composition, and
performance, developed by Ge Wang and Perry R. Cook~\cite{Wang_Cook_2003}.
It was designed to provide precise control over time, enabling composers to express complex audio processes with
deterministic timing and concurrency.

The language introduces a model where time is treated as a first-class concept, allowing programmers to explicitly
control the progression of time and the execution of musical events.
This enables accurate synchronization of concurrent processes and fine-grained control over temporal relationships in
music~\cite{Wang_Cook_2003}.

%% ---------------------------------------------------------------------------------------------------------------------

\subsubsection{Key Features}

\textbf{ChucK} follows a \textit{strongly-timed execution model}, where time advances only when explicitly specified by
the programmer.
The special keyword \texttt{now} represents the current time, and execution remains suspended until time is advanced,
ensuring deterministic behavior across runs~\cite{Wang_Cook_2003}.

A central feature of the language is its support for \textit{concurrency} through lightweight processes called
\textit{shreds}.
These shreds execute in parallel and are synchronized through the same timing mechanism, allowing multiple musical
processes to run concurrently with precise temporal alignment.

Example snippet:

\begin{lstlisting}[label={lst:chuck-1}]
  SinOsc osc => dac;
  440 => osc.freq;
  while (true) {
    1::second => now;
  }
\end{lstlisting}

Another defining characteristic is the \textit{ChucK operator} (\texttt{=>}), which unifies multiple language constructs
such as assignment, data flow, and signal routing.
This operator enforces a left-to-right execution model and allows chaining of operations, making the flow of audio and
data explicit~\cite{Wang_Cook_2003}.

The language also supports \textit{on-the-fly programming}, enabling code to be added, modified, or removed during
execution without interrupting the running system.
This feature makes \textbf{ChucK} particularly suitable for live coding and interactive performance
contexts~\cite{Wang_Cook_2003}.

Additionally, \textbf{ChucK} supports multiple simultaneous control rates, allowing different concurrent processes to
operate at varying temporal resolutions.

%% ---------------------------------------------------------------------------------------------------------------------

\subsubsection{Limitations}

Despite its precision and flexibility, \textbf{ChucK} presents several challenges.

First, the language operates at a relatively low level of abstraction, requiring programmers to explicitly manage time
and concurrency.
This can make it difficult to represent higher-level musical structures such as patterns, timelines, or hierarchical
compositions.

Second, the strongly-timed model places the responsibility of advancing time on the programmer, which can increase
complexity and reduce readability in larger programs.

Furthermore, while \textbf{ChucK} excels in real-time synthesis and performance, it is less focused on structured,
deterministic composition workflows, particularly those that require clear separation between composition and execution.

%% ---------------------------------------------------------------------------------------------------------------------

\subsubsection{Relevance to Our DSL}

\textbf{ChucK} highlights the importance of explicit time control and deterministic concurrency in music programming,
demonstrating how precise temporal semantics can be integrated directly into a programming language.

However, its low-level and time-driven approach contrasts with the goals of this thesis, which aim to provide a
higher-level, structured representation of musical concepts.
The proposed DSL seeks to retain deterministic behavior while offering more intuitive abstractions for composition, such
as timeline-based structures and modular musical components.
